\subsection{自同态的迹}

设 E 与 F 为两个希尔伯特空间。按照 V, p.~40 的约定，我们用 \(ba^{*}\)（对 E 中的 \(a\) 与 F 中的 \(b\)）表示从 E 到 F 中的连续线性映射 \(x \mapsto b \langle a|x\rangle\)。

引理 1. --- 存在一个同构 \(\theta\)，从向量空间 \(\mathbf{F} \otimes \mathbf{E}'\) 映到空间 \(\mathcal{L}_f(\mathbf{E}; \mathbf{F})\)（由从 \(\mathbf{E}\) 到 \(\mathbf{F}\) 中的所有有限秩连续线性映射组成）上，它由 \(\theta(b \otimes a^*) = ba^*\)（对 \(a \in \mathbf{E}\)、\(b \in \mathbf{F}\)）刻画。

由 A, II, § 4, No.~2，存在一个单射线性映射 \(\theta\)，从 \(\mathbf{F} \otimes \mathbf{E}'\) 到 \(\mathcal{L}(\mathbf{E};\mathbf{F})\) 中，且仅存在一个，它把 \(b \otimes a'\) 变为线性映射 \(x \mapsto ba'(x)\)（对 \(a' \in \mathbf{E}'\)、\(b \in \mathbf{F}\)）。显然 \(\theta(b \otimes a^*) = ba^*\)，且 \(\theta\) 的像包含在 \(\mathcal{L}_f(\mathbf{E};\mathbf{F})\) 中。然而，设 \(u \in \mathcal{L}_f(\mathbf{E};\mathbf{F})\)，且 \((e_1,\dots,e_n)\) 为 \(u\) 的像（在 \(\mathbf{F}\) 中）的标准正交基。令 \(f_i = u^*(e_i)\)（对 \(1 \leqslant i \leqslant n\)）。对每个 \(x \in \mathbf{E}\)，我们有

\[
u (x) = \sum_ {i = 1} ^ {n} \left\langle e _ {i} | u (x) \right\rangle . e _ {i} = \sum_ {i = 1} ^ {n} \left\langle f _ {i} | x \right\rangle . e _ {i},
\]

因而 \(u = \sum_{i=1}^{n} e_i f_i^* = \theta(\sum_{i=1}^{n} e_i \otimes f_i^*)\)。所以 \(\theta\) 的像等于 \(\mathcal{L}_f(\mathrm{E};\mathrm{F})\)。

\[
\mathrm{E} = \mathrm{F},
\]

\[
\mathscr {L} _ {f} (\mathrm{E}) = \mathscr {L} _ {f} (\mathrm{E}; \mathrm{E})
\]

\[
\mathcal {L} _ {f} (\mathrm{E})
\]

\[
\tau (\theta (a \otimes a ^ {\prime})) = a ^ {\prime} (a)
\]

\[
a \in \mathrm{E}, a ^ {\prime} \in \mathrm{E} ^ {\prime}\tag{22}
\]

\[
\tau (b a ^ {*}) = \langle a | b \rangle \quad \text { for } a, b \text { in } E.
\]

当 E 为有限维时，我们有 \(\mathcal{L}_{f}(\mathrm{E}) = \mathcal{L}(\mathrm{E})\)，而 \(\tau(u)\) 是 E 的自同态 u 的迹（A, II, § 4, No.~3）。

引理 2. --- 设 \((e_i)_{i\in \mathbf{I}}\) 为 E 的标准正交基。则

\[
\tau (u) = \sum_ {i \in \mathrm{I}} \left\langle e _ {i} | u (e _ {i}) \right\rangle
\]

对所有 \(u \in \mathcal{L}_f(\mathrm{E})\) 成立。

只需考虑 \(u = ba^{*}\)（其中 \(a, b\) 属于 E）的情形。这时

\[
\langle e _ {i} | u (e _ {i}) \rangle = e _ {i} ^ {*} b. a ^ {*} e _ {i} = \overline {{\langle e _ {i} | a \rangle}} \langle e _ {i} | b \rangle
\]

而引理 2 由公式 (22) 以及 V, p.~22 的公式 (3) 得出。

引理 3. --- 设 u 为 E 的连续正自同态，F 为 E 上所有有限秩正交投影算子组成的集。则对 E 的每个标准正交基 \((e_{i})_{i\in\mathbf{I}}\)，我们（在 \((in\overline{\mathbf{R}}_{+})\) 中）有等式

\[
\sum_ {i \in \mathrm{I}} \left\langle e _ {i} | u (e _ {i}) \right\rangle = \sup _ {p \in \mathcal {F}} \tau (p u p).
\]

对 I 的每个有限子集 J，令 \(p_{J} = \sum_{i \in J} e_{i} e_{i}^{*}\)；这是从 E 到由向量 \(e_{i}\)（i 取遍 J）生成的向量子空间上的正交投影算子。我们有

\[
p _ {\mathrm{J}} u p _ {\mathrm{J}} = \sum_ {i \in \mathrm{J}, j \in \mathrm{J}} \left\langle e _ {i} | u (e _ {j}) \right\rangle e _ {i} e _ {j} ^ {*},
\]

因而 \(\tau(p_{\mathbf{J}}up_{\mathbf{J}}) = \sum_{i\in \mathbf{J}}\langle e_i|u(e_i)\rangle\)。由于 \(p_{\mathbf{J}}\in \mathcal{F}\)，

\[
\sum_ {i \in \mathbf {J}} \left\langle e _ {i} | u (e _ {i}) \right\rangle \leqslant \sup _ {p \in \mathcal {F}} \tau (p u p);
\]

于是我们断定

\[
\sum_ {i \in \mathbf {I}} \left\langle e _ {i} | u (e _ {i}) \right\rangle = \sup _ {\mathbf {J}} \sum_ {i \in \mathbf {J}} \left\langle e _ {i} | u (e _ {i}) \right\rangle \leqslant \sup _ {p \in \mathcal {F}} \tau (p u p).
\]

设 \(v\) 为 \(\mathbf{E}\) 的一个有限秩连续正自同态，且 \(p \in \mathcal{F}\)。由 \(\mathbf{V}\), p.~23 的定理 2，存在标准正交基 \((f_{\alpha})_{\alpha \in \mathbf{A}}\)（属于 \(\mathbf{E}\)）及有限子集 \(\mathbf{B}\)（属于 \(\mathbf{A}\)），使得 \((f_{\alpha})_{\alpha \in \mathbf{B}}\) 是 \(p\) 的像的标准正交基。这时我们有 \(p = \sum_{\alpha \in \mathbf{B}} f_{\alpha} f_{\alpha}^{*}\)，于是与上面一样，有关系式 \(\tau(pvp) = \sum_{\alpha \in \mathbf{B}} \langle f_{\alpha}|v(f_{\alpha}) \rangle\)。由引理 2（V, p.~48）我们有 \(\tau(v) = \sum_{\alpha \in \mathbf{A}} \langle f_{\alpha}|v(f_{\alpha}) \rangle\)，由此得到公式

\[
\sum_ {\alpha \in \mathbf {B}} \langle f _ {\alpha} | v (f _ {\alpha}) \rangle \leqslant \tau (v).
\]

把这一不等式应用于 \(v = p_{\mathbf{J}}.up_{\mathbf{J}}\) 且 \(\mathbf{J}\) 为 \(\mathbf{I}\) 的有限子集的情形，就得到

\[
\sum_ {\alpha \in \mathbf {B}} \left\langle p _ {\mathbf {J}} (f _ {\alpha}) | u p _ {\mathbf {J}} (f _ {\alpha}) \right\rangle \leqslant \sum_ {i \in \mathbf {J}} \left\langle e _ {i} | u (e _ {i}) \right\rangle .\tag{23}
\]

对每个 \(x \in \mathbb{E}\)，我们有 \(p_{\mathbf{J}}(x) = \sum_{i \in \mathbf{J}} \langle e_i | x \rangle e_i\)，因而 \(x = \lim_{\mathbf{J}} p_{\mathbf{J}}(x)\)，这里极限是关于 \(\mathbf{J}\)（\(\mathbf{I}\) 的有限子集）组成的有序有向集取的。在 (23) 中对 \(\mathbf{J}\) 取极限，就得到

\[
\tau (p u p) = \sum_ {\alpha \in \mathrm{B}} \left\langle f _ {\alpha} | u (f _ {\alpha}) \right\rangle \leqslant \sum_ {i \in \mathrm{I}} \left\langle e _ {i} | u (e _ {i}) \right\rangle ,
\]

这样就完成了引理 3 的证明。

定义 7. --- 设 u 为希尔伯特空间 E 的连续正自同态。令

\[
\operatorname{Tr} (u) = \sup _ {p \in \mathcal {F}} \tau (p u p)\tag{24}
\]

（上界在 \(\overline{\mathbf{R}}_+\) 中），其中 \(\mathcal{F}\) 为 E 上所有有限秩正交投影算子组成的集。我们称 \(\operatorname{Tr}(u)\) 为 \(u\) 的迹。

设 \(p\) 为从 \(E\) 到 \(E\) 的某个有限维向量子空间上的正交投影算子，且 \((x_1, \ldots, x_m)\) 为 \(F\) 的标准正交基。我们已建立了关系式 \(\tau(pup) = \sum_{i=1}^{m} \langle x_i | u(x_i) \rangle\)。因此，我们可以用下列公式定义迹

\[
\operatorname{Tr} (u) = \sup _ {x _ {1}, \dots , x _ {m}} \sum_ {i = 1} ^ {m} \left\langle x _ {i} | u (x _ {i}) \right\rangle ,\tag{24'}
\]

其中 \((x_{1},\dots,x_{m})\) 取遍 E 的向量的所有有限标准正交序列组成的集。

由引理 3（V, p.~48），我们有

\[
\operatorname{Tr} (u) = \sum_ {i \in \mathbf {I}} \left\langle e _ {i} | u (e _ {i}) \right\rangle\tag{25}
\]

（对 E 的每个标准正交基 \((e_{i})_{i\in I}\)）。由此我们导出

\[
\operatorname{Tr} (u + v) = \operatorname{Tr} (u) + \operatorname{Tr} (v)\tag{26}
\]

\[
\operatorname{Tr} (\lambda u) = \lambda . \operatorname{Tr} (u)\tag{27}
\]

对 E 的所有连续正自同态 u、v 以及每个实数 \(\lambda \geqslant 0\) 成立（在 (27) 中我们约定 \(0. (+\infty) = 0\)）。设 \(\phi\) 为从 E 到希尔伯特空间 F 上的同构；由于 \(\phi\) 把 E 的每个标准正交基变为 F 的标准正交基，由 (25) 我们得到

\[
\operatorname{Tr} \left(\phi u \phi^ {- 1}\right) = \operatorname{Tr} (u).\tag{28}
\]

设 \((u_{\alpha})_{\alpha \in \mathbf{A}}\) 为 E 的连续正自同态组成的一个非空、递增有向且有界的族；令 \(u = \sup_{\alpha} u_{\alpha}\)，则 \(\langle x|u(x) \rangle = \sup_{\alpha} \langle x|u_{\alpha}(x) \rangle\) 对所有 \(x \in \mathrm{E}\) 成立（V, p.~46, 命题 13）。我们有 \(\operatorname{Tr}(u) = \sup_{\mathbf{J} \subset \mathbf{I}} \sum_{i \in \mathbf{J}} \langle e_i | u(e_i) \rangle\)，其中 \(\mathbf{J}\) 取遍 I 的所有有限子集，因而

\[
\operatorname{Tr} (u) = \sup _ {\alpha} \operatorname{Tr} (u _ {\alpha}) \quad \text { for } \quad u = \sup _ {\alpha} u _ {\alpha}.\tag{29}
\]

设 \(p_{\mathrm{F}}\) 为从 E 到希尔伯特子空间 F 上的正交投影算子；存在 E 的标准正交基 \((e_i)_{i \in \mathrm{I}}\) 及 I 的子集 J，使得 \((e_i)_{i \in \mathrm{J}}\) 是 F 的标准正交基。我们有 \(\operatorname{Tr}(p_{\mathrm{F}} up_{\mathrm{F}}) = \sum_{i \in \mathrm{J}} \langle e_i | u(e_i) \rangle\)。这个公式有两个推论：第一，我们有 \(\operatorname{Tr}(p_{\mathrm{F}} up_{\mathrm{F}}) \leqslant \operatorname{Tr}(u)\)；第二，取 \(u = 1_{\mathrm{E}}\)，就得到

\[
\operatorname{Tr} (p _ {\mathrm{F}}) = \left\{ \begin{array}{l l} \dim \mathrm{F} & \text { if   F   is   finite   dimensional } \\ + \infty & \text { if   not. } \end{array} \right.\tag{30}
\]

定义 8. --- 设 E 为复希尔伯特空间。用 \(\mathcal{L}^{1}(\mathrm{E})\) 表示 \(\mathcal{L}(\mathrm{E})\) 中由 E 的所有具有有限迹的连续正自同态生成的向量子空间。

由 V, p.~50 的公式 (25)，迹可以延拓为 \(\mathcal{L}^1 (\mathrm{E})\) 上的一个线性型，仍记作 Tr，且满足关系式 \(\operatorname{Tr}(u) = \sum_{i\in I}\langle e_i|u(e_i)\rangle\)（对所有 \(u\)（属于 \(\mathcal{L}^1 (\mathrm{E})\)）及 E 的每个标准正交基 \((e_i)_{i\in I}\)）。对每个 \(u\in \mathcal{L}^1 (\mathrm{E})\)，我们有 \(u^{*}\in \mathcal{L}^{1}(\mathrm{E})\) 且 \(\operatorname {Tr}(u^{*}) = \overline{\operatorname{Tr}(u)}\)。V, p.~50 的公式 (28) 可推广到 \(u\) 属于 \(\mathcal{L}^1 (\mathrm{E})\) 的情形。设 F 为 E 的希尔伯特子空间；由公式 (30)，正交投影算子 \(p_{\mathrm{F}}\) 属于 \(\mathcal{L}^1 (\mathrm{E})\) 当且仅当 F 是有限维的。对 E 中每对 \(a\) 与 \(b\)，我们有 \(4ab^{*} = \sum_{\varepsilon^{4} = 1}\varepsilon (a + \varepsilon b)(a + \varepsilon b)^{*}\)，且 \(cc^{*}\) 对所有 \(c\in \mathrm{E}\) 都是具有有限迹的正算子：因此，若 \(u\) 是 E 的有限秩连续自同态，则 \(u \in \mathcal{L}^1(\mathrm{E})\) 且 \(\operatorname{Tr}(u) = \tau(u)\)。

设 \(\mathbf{E}\) 为实希尔伯特空间，\(\mathrm{E}_{(\mathbf{C})}\) 为其复化（V, p.~5）。我们把 \(\mathbf{E}\) 等同于 \(\mathrm{E}_{(\mathbf{C})}\) 的一个子集。于是 \(\mathcal{L}(\mathrm{E})\) 可等同于 \(\mathcal{L}(\mathrm{E}_{(\mathbf{C})})\) 的一个实向量子空间，它由所有连续线性映射 \(u\)（从 \(\mathrm{E}_{(\mathbf{C})}\) 到 \(\mathrm{E}_{(\mathbf{C})}\) 中，满足 \(u(\mathrm{E}) \subset \mathrm{E}\)）组成。这时我们写 \(\mathcal{L}^1(\mathrm{E}) = \mathcal{L}(\mathrm{E}) \cap \mathcal{L}^1(\mathrm{E}_{(\mathbf{C})})\)。对每个 \(u \in \mathcal{L}^1(\mathrm{E})\)，迹 \(\operatorname{Tr}(u)\) 是实的且等于 \(\operatorname{Tr}(u^*)\)。公式 (25) 与 (28) 仍然有效，\(\mathcal{L}_f(\mathrm{E}) \subset \mathcal{L}^1(\mathrm{E})\) 且 \(\operatorname{Tr}(u) = \tau(u)\) 对所有 \(u \in \mathcal{L}_f(\mathrm{E})\) 成立。最后，闭向量子空间 \(\mathbf{F}\)（\(\mathbf{E}\) 的）是有限维的当且仅当 \(p_{\mathrm{F}}\) 属于 \(\mathcal{L}^1(\mathrm{E})\)。

* 注 1. --- 我们以后将定义从 Banach 空间 E 到 Banach 空间 F 中的核映射的概念。我们将证明，当 \(\mathcal{L}^1 (\mathrm{E})\) 由从 E 到 E 中的所有核映射组成时，E 是实或复的希尔伯特空间。 *

命题 14. --- 设 \(E_{1}, \ldots, E_{n}\) 为希尔伯特空间，\(E = E_{1} \hat{\otimes}_{2} \ldots \hat{\otimes}_{2} E_{n}\)，且 \(u_{i}\) 为 \(E_{i}\) 的连续自同态（对 \(1 \leqslant i \leqslant n\)）。若 \(u_{1}, \ldots, u_{n}\) 都是正的，则 \(u = u_{1} \hat{\otimes}_{2} \ldots \hat{\otimes}_{2} u_{n}\) 也是正的，并且

\[
\operatorname{Tr} (u) = \prod_ {i = 1} ^ {n} \operatorname{Tr} (u _ {i}).\tag{31}
\]

若 \(u_{i} \in \mathcal{L}^{1}(\mathrm{E}_{i})\)（对所有 \(1 \leqslant i \leqslant n\)），则 \(u \in \mathcal{L}^{1}(\mathrm{E})\)，且公式 (31) 在这一情形仍然有效。

对 \(n\) 用归纳法，我们立即归结到 \(n = 2\) 的情形。

对 i=1,2，我们定义半双线性型 \(\Phi_{i}\)（在 \(E_{i}\) 上），其公式为 \(\Phi_{i}(x,y)=\langle x|u_{i}(y)\rangle\)（对 \(E_{i}\) 中的 x、y）。若 \(u_{1}\) 与 \(u_{2}\) 都是正的，则型 \(\Phi_{1}\) 与 \(\Phi_{2}\) 是埃尔米特的且是正的。由 V, p.~25 的命题 1，存在一个正埃尔米特型 \(\Phi\)（在向量空间 \(E_{1}\otimes E_{2}\) 上），使得

\[
\Phi (x _ {1} \otimes x _ {2}, y _ {1} \otimes y _ {2}) = \Phi_ {1} (x _ {1}, y _ {1}). \Phi_ {2} (x _ {2}, y _ {2})
\]

（对 \(x_{1}, y_{1}\)（属于 \(E_{1}\)）与 \(x_{2}, y_{2}\)（属于 \(E_{2}\)））。我们立即验证关系式 \(\Phi(z, t) = \langle z|u(t) \rangle\)（其中 \(z\) 与 \(t\) 属于 \(E_{1} \otimes E_{2}\)）。由于 \(\Phi\) 是正的，我们有 \(\langle z|u(z) \rangle \geqslant 0\)（对所有 \(z\)，属于 \(E_{1} \otimes E_{2}\)）。由于 \(u\) 连续且 \(E_{1} \otimes E_{2}\) 在希尔伯特空间 \(E = E_{1} \hat{\otimes}_{2} E_{2}\) 中稠密，我们断定 \(u\) 是 \(E\) 的连续正自同态。

设 \((e_i)_{i\in I}\) 为 \(\mathbf{E}_1\) 的标准正交基，\((f_j)_{j\in J}\) 为 \(\mathbf{E}_2\) 的标准正交基；则族 \((e_i\otimes f_i)_{i\in I,j\in J}\) 是 E 的标准正交基，并且我们有

\[
\begin{array}{r l} & {\mathrm{Tr} (u) = \sum_ {i \in \mathbf {I}} \sum_ {j \in \mathbf {J}} \langle e _ {i} \otimes f _ {j} | u (e _ {i} \otimes f _ {j}) \rangle} \\ & {\quad = \sum_ {i \in \mathbf {I}} \sum_ {j \in \mathbf {J}} \langle e _ {i} | u _ {1} (e _ {i}) \rangle . \langle f _ {j} | u _ {2} (f _ {j}) \rangle} \\ & {\quad = \mathrm{Tr} (u _ {1}). \mathrm{Tr} (u _ {2}).} \end{array}
\]

特别地，若 \(u_{1}\) 与 \(u_{2}\) 是具有有限迹的正自同态，则 u 也是。由线性性，我们推出：当 K = C 时 u 属于 \(\mathcal{L}^{1}(\mathrm{E})\)，且对 i = 1, 2，诸 \(u_{i}\) 属于 \(\mathcal{L}^{1}(\mathrm{E}_{i})\)；公式 (31) 由线性性推广到这一情形。最后，K = R 且诸 \(u_{i} \in \mathcal{L}^{1}(\mathrm{E}_{i})\) 的情形可通过标量扩张归结为复的情形。

注 2. --- 设 E 为希尔伯特空间，它是希尔伯特子空间族 \((\mathrm{E}_{i})_{i\in\mathrm{I}}\) 的希尔伯特和。设 u 为 \(\mathcal{L}(\mathrm{E})\) 中满足 \(u(\mathrm{E}_{i}) \subset \mathrm{E}_{i}\)（对所有 \(i \in I\)）的元素；令 \(u_{i}\) 为 \(\mathcal{L}(\mathrm{E}_{i})\) 中在 \(E_{i}\) 上与 u 一致的元素。则 \(\operatorname{Tr}(u) = \sum_{i \in \mathrm{I}} \operatorname{Tr}(u_{i})\)（当 u 是正的，或属于 \(\mathcal{L}^{1}(\mathrm{E})\) 时）；这一关系式由 V, p.~50 的公式 (25) 应用于 E 的、由各 \(E_{i}\) 的标准正交基的并组成的标准正交基而得到。

